\chapter*{Resumen}
\addcontentsline{toc}{chapter}{Resumen}

Esta investigación, desarrollada en dos fases, construye un modelo estratégico de transporte para Bogotá y su región en el que la
demanda futura de viajes no es un pronóstico puntual sino un proceso estocástico. El índice de demanda
motorizada de cada una de doce zonas sigue un movimiento browniano geométrico (MBG) multivariado con un
factor común y choques idiosincráticos, simulado con su solución exacta y, como extensión, con saltos
sistémicos de tipo Merton que representan eventos como la pandemia de 2020. Para cada uno de
\NEscenarios{} escenarios y cada año entre 2026 y 2050 se resuelve el equilibrio entre transporte
público y vehículo privado en una red de dieciocho corredores, con congestión vial tipo BPR y
hacinamiento en el transporte público, calibrada para reproducir en 2025 la velocidad en hora pico
(18,9 km/h), la participación del transporte público y la duración media de los viajes observadas.

Un \emph{diseño urbano} combina proyectos férreos aún no contratados (extensión de la Línea 1 a la
calle 100, Línea 2, Línea 3 y Regiotram del Norte) con una política de proximidad y policentrismo de
intensidad $\theta$, inspirada en la «ciudad de 30 minutos» del POT. Los \NDisenos{} diseños se
comparan con números aleatorios comunes por su costo social en valor presente (tiempo, operación,
externalidades, inversión con sobrecostos lognormales y valor residual), y se elige el óptimo
minimizando una combinación de la media y el CVaR al 95\,\%.

Con los supuestos centrales y la tasa social de descuento del DNP (9\,\%), el diseño óptimo es
\OptimoNeutral: reduce el costo social esperado de \VPBase{} a \VPOptimo{} billones de pesos de 2025,
con un beneficio neto esperado de \BNOptimo{} billones y positivo en el \PBNOptimo{} de los escenarios;
el resultado se mantiene para todo peso del CVaR $w<\WCambio$; desde ese valor el óptimo agrega la
extensión de la Línea 1, con un CVaR casi idéntico. Las líneas de metro
adicionales reducen la dispersión del costo, pero su costo de inversión supera los ahorros esperados a
esa tasa; su signo cambia con tasas de descuento más bajas y, en presencia de saltos sistémicos, la
combinación con la Línea 3 pasa a ser la de menor CVaR. La valoración de la Línea 3 como opción real
con mínimos cuadrados de Monte Carlo, validada contra la solución analítica de McDonald y Siegel,
indica que comprometerse hoy destruye valor (VPN esperado de \OpcionAhora{} billones) y que la regla
óptima es esperar y construir solo si el ahorro operativo observado cruza una frontera de ejercicio.

La fase I recalibra el modelo a la coyuntura de 2013, cuando se diseñó un metro subterráneo, y evalúa
la Primera Línea con la información de entonces. Con la tasa de descuento de 12\,\% vigente, el metro
elevado tenía el mayor beneficio neto esperado (\BNElevDoce{} billones, frente a \BNSubOfDoce{} del
subterráneo con su costo oficial y \BNSubUnDoce{} con costos de mercado). La conclusión no es robusta:
con la tasa actual de 9\,\% y el costo oficial ambas tipologías empatan, y por debajo de \TasaEmpate{}
el subterráneo sería preferible. El tamaño de la ciudad, por sí solo, no justifica el subterráneo: la
capacidad del viaducto no es la restricción y una red de varias líneas favorece la tipología más
barata. El argumento decisivo es el costo urbano del viaducto: con 12\,\% tendría que superar
\UmbralOfDoce{} billones por año, frente a una estimación por pérdida de valor de los predios de 0,06 a
0,28 billones. La demanda observada en 2025 quedó en el percentil \PercentilObsFin{} de los escenarios
de 2013, y con ella el subterráneo habría perdido más que el elevado. Para las líneas 2 y 3, la
tipología subterránea cuesta entre \AhorroLdos{} y \AhorroLtres{} billones más y exige una
justificación explícita.

\medskip
\noindent\textbf{Palabras clave:} movimiento browniano geométrico, Monte Carlo, cuasi-Monte Carlo,
CVaR, opciones reales, mínimos cuadrados de Monte Carlo, equilibrio de transporte, diseño urbano,
metro elevado, metro subterráneo, Bogotá.

\chapter*{Abstract}
\addcontentsline{toc}{chapter}{Abstract}

This study builds a strategic transportation model for Bogotá in which future travel demand is a
stochastic process rather than a point forecast. The motorized demand index of twelve zones follows a
multivariate geometric Brownian motion (GBM) with a common factor, simulated exactly and extended with
Merton-type systemic jumps. For \NEscenarios{} scenarios and every year from 2026 to 2050, a
mode-choice equilibrium with BPR road congestion and transit crowding is solved on an eighteen-corridor
network calibrated to 2025 peak speed, transit share and mean trip duration. Forty-eight urban designs,
combining uncommitted rail projects with a proximity and polycentric land-use policy, are compared with
common random numbers by their social cost and ranked by a mean--CVaR criterion. Under central
assumptions and the official 9\,\% social discount rate, the optimal design is the Northern Regiotram
combined with the full land-use policy, with an expected net benefit of \BNOptimo{} trillion pesos and
a \PBNOptimo{} probability of a positive outcome. Additional metro lines reduce cost dispersion but do
not pay off at that discount rate; a least-squares Monte Carlo real-options analysis, validated against
the McDonald--Siegel closed form, shows that the value-maximizing policy for Metro Line 3 is to wait and
build only if observed savings cross an exercise boundary. A first phase recalibrates the model to
2013, when an underground metro was designed, and evaluates Metro Line 1 with the information then
available. At the 12\,\% discount rate in force, the elevated line had the highest expected net
benefit; at today's 9\,\% rate and the official 2014 cost both typologies tie, and below
\TasaEmpate{} the underground line would be preferable. City size alone does not justify tunnelling:
the decisive factor is the urban cost of the viaduct, which would have to exceed \UmbralOfDoce{}
trillion pesos per year at 12\,\%, against a property-value estimate of 0.06--0.28 trillion. Observed
2025 demand fell at the \PercentilObsFin{} percentile of the 2013 scenarios, and under it the
underground line would have lost more than the elevated one.

\medskip
\noindent\textbf{Keywords:} geometric Brownian motion, Monte Carlo, quasi-Monte Carlo, CVaR, real
options, least-squares Monte Carlo, transport equilibrium, urban design, elevated metro, underground metro, Bogotá.
